\begin{frame}
    \frametitle{Um problema não trivial}
    \metroset{block=fill}

    \begin{block}{Problema}
        De quantas maneiras podemos preencher um tabuleiro $1 \times n$ com peças $1 \times 1$ e $1 \times 2$?
    \end{block} \pause

    O número de maneiras pode ser modelado recursivamente como $f(n) = f(n-1) + f(n-2)$. Vejamos uma abordagem diferente: \pause
    \begin{align*}
        \begin{pmatrix}
            1 & 1 \\
            1 & 0
        \end{pmatrix} \cdot 
        \begin{pmatrix}
            f(n) \\
            f(n-1)
        \end{pmatrix} = \pause
        \begin{pmatrix}
            f(n) + f(n-1) = f(n+1) \\
            f(n)
        \end{pmatrix}
    \end{align*} \pause

    Seja $A$ a matriz acima e $v_n = (f(n), f(n-1))$. \pause A equação acima diz que $Av_n = v_n+1$. Como $A v_{n+1} = v_{n+2}$, Temos
        \[ A(Av_n) = v_{n+2} \ . \] \pause
    Pela associatividade da multiplicação de matrizes, $A^2 v_n = v_{n+2}$. Assim, $A^{10^{18} - 1} v_1 = v_{10^{18}}$. \pause

    {\bf Solução.} Use o algoritmo de exponenciação rápida para calcular $A^n$ em $O(\log n)$.
\end{frame}